% !PRINTED 327
\weberpage{327}
If the point $\xi_3$ is a point of inflection and $x_1 = 0$ is the inflectional tangent, then $f$ has the form

\weq{4}{f = x_1 \varphi + x_2^3 \psi,}
in which $\varphi$ is a form of degree $(n-1)$, $\psi$ a form of degree $(n-3)$.

Conversely, if $f$ has this form, the point $\xi_3$ is a point of inflection, and $x_1 = 0$ is the inflectional tangent.

A straight line that touches the curve $f$ in two separate points is called a \gesperrt{double tangent}. If we take such a double tangent as the line $x_3 = 0$ and the points of contact as the vertices of the coordinate triangle lying on $x_3 = 0$, then, when $x_3 = 0$ is substituted into $f = 0$, an equation must arise that has both $x_1$ and $x_2$ as quadratic factors. The function $f$ must therefore have the form:

\weq{4}{f = x_3 \varphi + x_1^2 x_2^2 \psi,} % [sic] number (4) reprinted; page also numbers preceding display (4)
in which $\varphi$ is a form of degree $(n-1)$, $\psi$ a form of degree $(n-4)$.

Somewhat more generally we may also say: if $x_3 = 0$ is a double tangent, then $f$ can be represented thus:

\weq{5}{f = x_3 \varphi + \chi^2 \psi,}
in which $\chi = 0$ is the equation of a conic section. The points of contact of the double tangent are the points of intersection of $x_3 = 0$ and $\chi = 0$.

\begin{center}
\S.\,87.

\gesperrt{Fundamental Covariants of a Ternary Form.}
\end{center}

In \S.\,60 of the first volume we became acquainted with the concept of the invariants and covariants of a form. In geometry too these forms play a great role.

If by the linear substitution \S.\,85, (1) the ternary form $f(x_1, x_2, x_3)$, which we also denote in abbreviated fashion by $f(x)$, passes into $\varphi(y)$, and if we denote by $a$ the coefficients of $f$, by $b$ the coefficients of $\varphi$, and by $r$ the substitution determinant, then a form $\Phi(x, a)$ that is homogeneous both with respect to the $x$ and with respect to the $a$ is called a \gesperrt{covariant} of $f(x)$ if it satisfies the condition

\weq{1}{\Phi(y, b) = r^\lambda \Phi(x, a)}
%CONT>>
